%! Departures from Linearity — Unit 2, Lesson 2.5 — Guided Notes & Practice (Answer Key)
% MIRROR of ../notes/main.tex: -key in place of -boxes, \termblank -> \termans, \blank -> \ans, and the
% fourth argument of every \pcell filled. Nothing else differs — keep it that way.
% THE in-class document, in the MAIN IDEAS / NOTES shape (see LESSON_SHAPE.md §1), at the
% COMPACT budget: the vocabulary box, then ONE guidednotes table of THREE rows (two instruction
% rows + Guided Practice), 13 problems, 2 pages.
%   I DO   : the teacher fills the definition lines and works the FIRST problem of each grid
%            (1, 5, 8 — 8 opens the "Straighten the curve" sub-row).
%   WE DO  : the class works the rest; the last row, Guided Practice, is worked together.
%   YOU DO : nothing on this page — ../homework/ IS the individual practice.
% THE TRAP is problem 3 (an outlier far off the trend need not move the slope).
% THE CRUX is problem 7, in the last instruction row: "r = -0.94 is close to -1, so the line is
% a good model" — while the residual plot is a clear U (EK DAT-1.F.2; Topic 2.9 DAT-1.J.2).
%
% THE DATA
%  Row 1 — ten coffee shops, seats (x) vs Saturday customers (y):
%    x: 20 24 28 30 35 38 42 45 50 55   y: 160 170 205 190 240 225 265 250 300 295
%    10 shops: yhat = 77.9 + 4.14x, r = 0.97.   + A(36,300): 86.0 + 4.10x, r = 0.87
%    + B(100,490): 78.9 + 4.12x, r = 0.99.      + C(100,230): 187.1 + 1.01x, r = 0.48
%  Row 2 — antibiotic concentration (mg/L) at hours 0..10:
%    20.2 14.9 11.7 9.0 6.5 4.9 3.9 2.9 2.1 1.7 1.2
%    linear: yhat = 15.9 - 1.74x, r = -0.94; residuals 4.3 .8 -.7 -1.7 -2.4 -2.3 -1.5 -.8 .1 1.5 2.7
%    log10(conc): 1.31 - 0.12x, r = -0.9996; residuals (x1000) 0 -11 5 13 -7 -9 13 6 -13 17 -13
%  Row 3 (GP) — 13 days, high temperature (F) vs cups of lemonade:
%    T: 72 75 77 78 80 82 83 85 87 88 90 93 95   L: 70 82 71 90 85 128 96 92 110 98 116 112 125
%    yhat = -81.6 + 2.15x, r = 0.80; without the 82F fair day: -91.7 + 2.24x, r = 0.94
%
% PAGE LOCKSTEP with notes_key/: \termblank -> \termans, \blank -> \ans, 4th arg of \pcell.
% TABLE MECHANICS: inside a cell `\\` ENDS THE ROW — use \par. Long rows are split into
% sub-rows (a bare `\\` then `& ...`) so the table can break between them.
% NO SKETCH QUESTIONS: every display is pre-drawn in TikZ and read.
\documentclass[12pt]{article}
\usepackage{apstats-article}
\usepackage{apstats-key}   % pulls in -boxes; do NOT also load -boxes

\begin{document}

\pageheader{Unit 2, Lesson 2.5}{Guided Notes \& Practice --- Answer Key}
% NAMESTRIP: no \namedateperiod here — the name row lives on the cover only.

\vspace{2pt}

\boxguard[16]
\begin{vocabbox}
\small
Fill in each term \textbf{as we name it} in the notes below.
\par
\termans{Outlier (in regression)}{a point that does not follow the trend of the rest; it has a large residual.}
\par
\termans{High-leverage point}{a point whose $x$-value is much larger or smaller than the other $x$-values.}
\par
\termans{Influential point}{a point that, if removed, changes the slope, $y$-intercept, or $r$ a lot.}
\par
\end{vocabbox}

\small
\begin{guidednotes}
% ── ROW 1: outliers, high leverage, influence ────────────────────────────────
\mainidea[Three kinds of]{Unusual points} &
An \textbf{outlier} does not follow the trend: it has a large \ans{residual}. A
\textbf{high-leverage point} has an \ans{$x$-value} far larger or smaller than the rest. An
\textbf{influential point} is one that, if \ans{removed}, changes the slope, the $y$-intercept,
or $r$ a lot. Ten coffee shops: seats ($x$) and Saturday customers ($y$). Shops A, B, and C are each added to
the ten, \emph{one at a time}, and the line is refit.

\vspace{3pt}
\begin{minipage}[c]{0.46\linewidth}
\centering
\begin{tikzpicture}[xscale=0.041, yscale=0.0052]
  \foreach \y in {100,200,...,500} {
    \draw[linegray, very thin] (0,\y) -- (110,\y);
    \node[left, font=\tiny] at (0,\y) {\y};}
  \draw[->, thick] (0,100) -- (0,560);
  \draw[->, thick] (0,100) -- (112,100);
  \foreach \x in {0,20,...,100} {
    \draw (\x,95) -- (\x,105); \node[below, font=\tiny] at (\x,97) {\x};}
  \draw[navy, dashed] (15,140) -- (105,512.6);
  \foreach \x/\y in {20/160, 24/170, 28/205, 30/190, 35/240, 38/225, 42/265, 45/250, 50/300, 55/295}
    \fill[navy] (\x,\y) circle (1.25 and 9.4);
  \foreach \x/\y/\l in {36/300/A, 100/490/B, 100/230/C} {
    \fill[redacc] (\x-1.6,\y-12) rectangle (\x+1.6,\y+12);
    \node[right, font=\scriptsize\bfseries, text=redacc] at (\x+1,\y) {\l};}
  \node[font=\tiny] at (55,5) {Seats};
  \node[rotate=90, font=\tiny] at (-17,330) {Customers};
\end{tikzpicture}
\end{minipage}\hfill
\begin{minipage}[c]{0.52\linewidth}
\centering\footnotesize
\setlength{\tabcolsep}{3pt}%
\renewcommand{\arraystretch}{1.2}
\begin{tabular}{@{} l l c @{}}
\rowcolor{sky}
\textbf{Line fit to} & \textbf{Equation} & $\boldsymbol{r}$ \\
\hline
10 shops & $\hat{y} = 77.9 + 4.14x$ & 0.97 \\
10 + A & $\hat{y} = 86.0 + 4.10x$ & 0.87 \\
10 + B & $\hat{y} = 78.9 + 4.12x$ & 0.99 \\
10 + C & $\hat{y} = 187.1 + 1.01x$ & 0.48 \\
\end{tabular}\\[3pt]
{\scriptsize Dashed line: the fit to the 10 shops.}
\end{minipage}
\\
& \notesprompt{Use the plot and the table.}
\begin{probgrid}{|Y|Y|}
\pcell{1}{Which added shop is an outlier but \emph{not} high-leverage? Why?}{1.3cm}{A: 36 seats is typical, but it sits far above the trend.} &
\pcell{2}{Which added shops are high-leverage? Why?}{1.3cm}{B and C: 100 seats, far beyond the others' 20--55.} \\ \hline
\pcell{3}{A sits far off the trend. Does adding A move the slope much?}{1.3cm}{No: slope $4.14 \to 4.10$. Only $r$ drops ($0.97 \to 0.87$).} &
\pcell{4}{Which shop is influential? Justify it with numbers from the table.}{1.3cm}{C: slope $4.14 \to 1.01$ and $r$ $0.97 \to 0.48$.} \\ \hline
\end{probgrid}
\\ \hline
% ── ROW 2: is a line appropriate? — THE CRUX ─────────────────────────────────
\mainidea[Is a line right?]{Residual plots} &
\textbf{Random scatter} in a residual plot says a line fits; a \ans{curved} pattern says a
line is the wrong model. An antibiotic's level in a patient's blood ($y$, mg/L), each hour ($x$)
after one dose: \ $\hat{y} = 15.9 - 1.74x$, \ $r = -0.94$.
\begin{center}
\begin{tikzpicture}[xscale=0.40, yscale=0.065]
  \draw[->, thick] (0,0) -- (10.6,0);
  \draw[->, thick] (0,0) -- (0,22);
  \foreach \x in {0,2,...,10} {
    \draw (\x,-0.5) -- (\x,0.5); \node[below, font=\tiny] at (\x,-0.3) {\x};}
  \foreach \y in {0,10,20} { \node[left, font=\tiny] at (0,\y) {\y}; }
  \draw[navy, thick] (0,15.9) -- (9.14,0);
  \foreach \x/\y in {0/20.2,1/14.9,2/11.7,3/9.0,4/6.5,5/4.9,6/3.9,7/2.9,8/2.1,9/1.7,10/1.2}
    \fill[navy] (\x,\y) circle (0.14 and 0.7);
  \node[font=\tiny] at (5.3,-6.5) {Hours after the dose};
  \node[rotate=90, font=\tiny] at (-1.9,11) {mg/L};
  \node[font=\bfseries\scriptsize] at (5.3,25) {Scatterplot and line};
  \begin{scope}[xshift=14cm, yscale=2.9]
    \draw[linegray] (0,0) -- (10.6,0);
    \draw[->, thick] (0,-3) -- (0,5);
    \draw[thick] (0,-3) -- (10.6,-3);
    \foreach \y in {-2,0,2,4} { \node[left, font=\tiny] at (0,\y) {\y}; }
    \foreach \x in {0,2,...,10} { \node[below, font=\tiny] at (\x,-3) {\x}; }
    \foreach \x/\y in {0/4.3,1/0.8,2/-0.7,3/-1.7,4/-2.4,5/-2.3,6/-1.5,7/-0.8,8/0.1,9/1.5,10/2.7}
      \fill[navy] (\x,\y) circle (0.14 and 0.23);
    \node[font=\tiny] at (5.3,-5.2) {Hours after the dose};
    \node[rotate=90, font=\tiny] at (-1.9,1) {Residual};
    \node[font=\bfseries\scriptsize] at (5.3,5.9) {Residual plot};
  \end{scope}
\end{tikzpicture}
\end{center}
\\
& \notesprompt{Read the residual plot.}
\begin{probgrid}{|Y|Y|Y|}
\pcell{5}{Describe the pattern in the residuals.}{1.4cm}{A U-shape: $+$ at both ends, $-$ in the middle.} &
\pcell{6}{At 0 and 10 hours, does the line predict too high or too low?}{1.4cm}{Too low --- the residuals there are positive.} &
\pcell{7}{``$r = -0.94$ is close to $-1$, so the line is a good model.'' Agree? Why?}{1.4cm}{No --- the residuals curve, so a line is the wrong form.} \\ \hline
\end{probgrid}
\\
& \fcolorbox{navy}{sky}{\parbox{\dimexpr\linewidth-2\fboxsep-2\fboxrule\relax}{\centering
\textbf{$r$ measures strength. The residual plot checks \ans{form}.}}}

\vspace{3pt}
A \textbf{transformation} can straighten a curve. Take the \ans{log} of each level and fit a
line to $\log y$: \ $\widehat{\log y} = 1.31 - 0.12x$, \ $r = -0.9996$, and its residual plot
is random scatter.
\notesprompt{Straighten the curve.}
\begin{probgrid}{|Y|Y|}
\pcell{8}{Give two pieces of evidence that the log model fits better than the line.}{1.0cm}{Its residuals are random, and $r = -0.9996$ is closer to $-1$.} &
\pcell{9}{Use the log model to predict the level 5 hours after the dose.}{1.0cm}{$1.31 - 0.12(5) = 0.71$; \ $10^{0.71} \approx 5.1$ mg/L.} \\ \hline
\end{probgrid}
\\ \hline
% ── ROW 3: guided practice — WE DO, together, fresh context ──────────────────
\mainidea[Guided practice]{The Lemonade Stand} &
\begin{minipage}[c]{0.46\linewidth}\raggedright
A lemonade stand recorded the high temperature (\textdegree F) and cups sold on 13 summer
days.\par
\vspace{2pt}
$\hat{y} = -81.6 + 2.15x$, \ $r = 0.80$
\end{minipage}\hfill
\begin{minipage}[c]{0.52\linewidth}
\centering
\begin{tikzpicture}[xscale=0.17, yscale=0.032]
  \draw[linegray] (70,0) -- (97,0);
  \draw[->, thick] (70,-18) -- (70,40);
  \foreach \y in {-15,0,15,30} { \node[left, font=\tiny] at (70,\y) {\y}; }
  \foreach \x in {70,75,...,95} { \node[below, font=\tiny] at (\x,-18) {\x}; }
  \foreach \x/\y in {72/-3.4,75/2.1,77/-13.2,78/3.7,80/-5.6,83/-1.1,85/-9.4,87/4.3,88/-9.8,90/3.8,93/-6.6,95/2.1}
    \fill[navy] (\x,\y) circle (0.25 and 1.35);
  \fill[redacc] (82,33.1) circle (0.3 and 1.6);
  \node[right, font=\tiny, text=redacc] at (82.5,33.1) {street fair};
  \node[font=\tiny] at (83.5,-35) {High temperature (\textdegree F)};
  \node[rotate=90, font=\tiny] at (63.5,11) {Residual (cups)};
\end{tikzpicture}
\end{minipage}
\\
& \notesprompt{We work these together.}
\begin{probgrid}{|Y|Y|}
\pcell{10}{Describe the residual plot.}{1.0cm}{Random scatter, no curve; one big $+$ residual at 82\textdegree F.} &
\pcell{11}{$r$ is only 0.80. Is a linear model appropriate here? Why?}{1.0cm}{Yes --- no curve, so a line fits. It is just moderate.} \\ \hline
\pcell{12}{The street-fair day: outlier, high-leverage, or both? Why?}{1.0cm}{Outlier only: 82\textdegree F is a typical $x$, but its residual is $+33$.} &
\pcell{13}{Without that day, $\hat{y} = -91.7 + 2.24x$ and $r = 0.94$. What changed
most?}{1.0cm}{$r$ ($0.80$ vs.\ $0.94$); the slope barely moved ($2.15$ vs.\ $2.24$).} \\ \hline
\end{probgrid}
\\ \hline
\end{guidednotes}

% ── THE NOTES END AT GUIDED PRACTICE ─────────────────────────────────────────
% There is deliberately no "you do" here: ../homework/ is the individual practice,
% started in class in the last 8 minutes while the teacher circulates.

\end{document}
